% =====================================================================
%  c07.tex —— §8 结语（原文末节）+ 参考文献（原书 p.57–58；PDF p.16–p.17）
%  编号公式：无（全文公式计数器止于 (10)，在 c06）
%  无编号展示：无
%  定理类环境：无
%  注：参考文献 22 条照抄原文不译，保留 MR/arXiv/math.QA 等分类号；
%      [9][10][11] 为日文文献，标题照抄日文原样。
%      [10] 标题中的中点（U+30FB）在本机 SimSun 中缺字，故就地切 MS Gothic
%      渲染这一个字符（不影响字符本身，只换字形）。
%      跨块引用一律写原书字面编号（例 6.1、[12]），以免单块编译出现
%      undefined reference / citation；文献引用不用 \cite。
% =====================================================================

% ===== 原书 p.57 =====
\section{结语}

讨论同调镜像对称时，（就个人而言）\(\ainf\)范畴的等价似乎更为本质
（例如 [15, 8]），三角范畴的等价则似乎只是由它导出的结果。实际上，
正如本文所讨论的，从\(\ainf\)范畴\(Tw(C)\)落到三角范畴\(Tr(C)\)，
一般会丢失信息。因此，\(Tw(C)\)或许才是应当考虑的最基本的东西。
然而，\(Tw(C)\)一般比\(C\)的对象数庞大得多。与\(Tr(C)\)相比，它还保留着
高次积结构，态射也包含所有次数的。至于次数，\(H(Tw(C))\)可由
\(H^{0}(Tw(C))=Tr(C)\)复原；即便舍弃高次积、化归为三角范畴\(Tr(C)\)，
其中也含有\(\ainf\)范畴\(Tw(C)\)相当多的信息（cf.\ [12]）。
所以\(Tw(C)\)的定式化带着信息重复、冗余过多的缺陷。
既然如此，或许应当寻找另一种定式化，找出与\(Tw(C)\)持有等价信息的东西。

\transnote{原文只断言「次数方面 \(H(Tw(C))\) 可由 \(H^{0}(Tw(C))=Tr(C)\) 复原」，
未给做法。用语是平移：\(Tr(C)\) 的平移函子 \([1]\) 同时作用于
\(H(Tw(C))\) 的次数指标，故
\(H^{i}(Tw(C))(X,Y)\cong H^{0}(Tw(C))(X,Y[i])\)。
因此 \(Tr(C)\) 连同平移结构，已经决定了 \(H(Tw(C))\) 的全部分次。}

另一方面，本次在例 6.1 中构造的、\(\ainf\)增广不唯一的三角范畴，
是否能作为例如某个流形上相干层的导出范畴，或矩阵因子化构成的三角范畴
而得到，目前还不清楚。
例 6.1 的\(C_{\epsilon}\)是带次序的\(\ainf\)范畴（ordered \(\ainf\) category），
即\(C_{\epsilon}\)构成\(Tr(C_{\epsilon})\)的例外生成系。
因此，\(Tr(C_{\epsilon})\)具有 Serre 函子（[1], cf.\ [11]），
在这个意义上它是性质相当好的对象。
不过，哪一类三角范畴由怎样的几何得到？我认为这也是一个有趣的问题。
也就是说，固定从各自几何确定三角范畴的方法之后，\(\ainf\)增广唯一的
几何有多少？若\(\ainf\)增广不唯一的几何反倒「足够例外」，
则不妨先把它们当作例外的，而对非例外的几何只考虑三角范畴、
不考虑\(\ainf\)范畴——这个方向或许是妥当的。

\transnote{「十分例外的」加了引号，是双关。前文的「例外生成系」是技术术语
（见 [1]、[11]）：一族对象按次序排列，其间的 \(\Ext\) 只在零次非零，
由此 \(Tr(C_{\epsilon})\) 带上 Serre 函子。此处又借「例外」的日常义，
指「罕见、属少数情形」。作者的问题是：若增广不唯一的几何只是少数例外，
就可把它们归入「例外」一类，其余几何只研究三角范畴。}

\begin{thebibliography}{22}

\bibitem{b1} A. Bondal, Representations of associative algebras and coherent
sheaves. (Russian) \emph{Izv. Akad. Nauk SSSR Ser. Mat.} 53:25--44, 1989;
translation in \emph{Math. USSR-Izv.} 34:23--42, 1990;
\emph{Helices, representations of quivers and Koszul algebras}.
\emph{Helices and vector bundles}, 75--95, London Math. Soc. Lecture Note
Ser., 148, Cambridge Univ. Press, Cambridge, 1990.

\bibitem{b2} A. Bondal and M. Kapranov, Enhanced triangulated categories.
\emph{Math. USSR Sbornik} 70:93--107, 1991.

\bibitem{b3} K. Fukaya. Morse homotopy, \(A^{\infty}\)-category, and Floer
homologies. In \emph{Proceedings of GARC Workshop on Geometry and Topology
'93 (Seoul, 1993)}, volume 18 of \emph{Lecture Notes Ser.}, pages 1--102.
Seoul Nat. Univ., Seoul, 1993.

\bibitem{b4} K. Fukaya. Floer homology and mirror symmetry. II. In
\emph{Minimal surfaces, geometric analysis and symplectic geometry},
volume 34 of \emph{Adv. Stud. Pure Math.}, pages 31--127. Math. Soc. Japan,
Tokyo, 2002. (Baltimore, MD, 1999).

\bibitem{b5} I. Gelfand and Y. Manin. \emph{Methods of Homological Algebra}.
Springer-Verlag, 1994.

% ===== 原书 p.58 =====
\bibitem{b6} T. V. Kadeishvili. The algebraic structure in the homology of an
\(A(\infty)\)-algebra. (Russian) \emph{Soobshch. Akad. Nauk Gruzin. SSR}
108:249--252 (1982).

\bibitem{b7} H. Kajiura. Noncommutative homotopy algebras associated with open
strings. \emph{Reviews in Math. Phys.}, 19(1):1--99, 2007. math.QA/0306332.

\bibitem{b8} H. Kajiura. Homological perturbation theory and homological mirror
symmetry. In \emph{Higher Structures in Geometry and Physics}, volume 287 of
\emph{Progress in Math.}, page 201--226, Birkhäuser/Springer, New York, 2011.

\bibitem{b9} 梶浦宏成. \(A_{\infty}\) 圏と有限次元代数の表現論.
第 11 回 代数群と量子群の表現論 研究集会報告集, 2008.

\bibitem{b10} 梶浦宏成. マッセイ積の一般化と三角圏.
空間の代数的{\CJKfontspec{MS Gothic}・}幾何的モデルとその周辺 2009 （於 信州大学）報告集.

\bibitem{b11} 梶浦宏成. 数物系のための圏論\textendash{} 導来圏、三角圏,
\(A_{\infty}\) 圏を中心に. SGC ライブラリ 75, 数理科学, 2010.

\bibitem{b12} H. Kajiura. On \(A_{\infty}\)-enhancements for triangulated
categories. \emph{J. Pure Appl. algebra.} 217: 1476--1603, 2013.

\bibitem{b13} M. Kashiwara and P. Shapira. \emph{Sheaves on Manifolds}.
Grundleheren 292, Springer-Verlag, 1990.

\bibitem{b14} M. Kontsevich. Homological algebra of mirror symmetry. In
\emph{Proceedings of the International Congress of Mathematicians,
(Zürich, 1994)}, volume 184, pages 120--139. Birkhäuser, 1995.
math.AG/9411018.

\bibitem{b15} M. Kontsevich and Y. Soibelman. Homological mirror symmetry and
torus fibrations. In \emph{Symplectic geometry and mirror symmetry
(Seoul, 2000)}, pages 203--263. World Sci. Publishing, River Edge, NJ, 2001.
math.SG/0011041.

\bibitem{b16} V. A. Lunts and D.O. Orlov. Uniqueness of enhancement for
triangulated categories. arXiv:0908.4187, 2009.

\bibitem{b17} J. P. May, Matric Massey products. \emph{J. Algebra}
12:533--568, 1969.

\bibitem{b18} D. Maes, private communication, November 8, 2013.

\bibitem{b19} F. Muro, S. Schwede, and N. Strickland, Triangulated categories
without models. \emph{Invent. Math.} 170: 231--241, 2007.

\bibitem{b20} P. Seidel, \emph{Fukaya categories and Picard-Lefschetz theory},
Zurich Lectures in Advanced Mathematics, European Mathematical Society (EMS),
Zurich, 2008.

\bibitem{b21} J. Stasheff. Homotopy associativity of H-spaces, I.
\emph{Trans. Amer. Math. Soc.}, 108:293--312, 1963.

\bibitem{b22} J. Stasheff. Homotopy associativity of H-spaces, II.
\emph{Trans. Amer. Math. Soc.}, 108:313--327, 1963.

\end{thebibliography}

\vspace{2.5em}
\begin{center}
\textsc{Faculty of Science, Chiba University, Chiba 263-8522, Japan}\\[6pt]
\emph{E-mail address}: [邮箱隐去]
\end{center}
